\section{研究问题}

本研究以雨天扰动为主要情景，关注西安轨道交通系统在外部天气冲击下的服务功能保持与韧性变化。研究重点并不是简单描述雨天客流是否下降，而是通过构建同期匹配的正常天气反事实基线，识别雨天扰动下站点、OD 联系、区间和换乘节点的响应差异，并进一步解释这种差异背后的站点环境因素。核心研究问题如下：

\begin{enumerate}
\item 在正常天气条件下，西安轨道交通系统中站点、OD 联系、区间和换乘节点的基线服务功能与运行负荷结构呈现何种特征？哪些空间单元构成潜在的韧性薄弱环节？
\item 相对于同期匹配的正常天气反事实基线，雨天扰动会使西安轨道交通系统的站点客流、OD 联系强度以及区间和换乘节点运行压力产生怎样的响应与重分配？
\item 雨天扰动下轨道交通系统的服务功能保持率与脆弱性在不同空间单元之间呈现何种差异？站点周边 POI 功能结构与建成环境特征如何解释或调节这种韧性异质性？
\end{enumerate}
